\section{Local strictness and an exact multiplier ideal}
\label{inj:localized}

The rank-one direct image of an adjoint line need not be locally free.
We identify its defect from its reflexive hull with the multiplier ideal
of its fiber-integral metric. This exact identity permits Euler
interpolation without discarding conditions in codimension two.

The two analytic lemmas below turn positivity on one horizontal patch
into vanishing for that ideal. Proposition~\ref{inj:exact-adjoint-ideal}
then identifies the ideal, and Proposition~\ref{inj:degree-zero-route}
applies the construction to the curvature-null-pole field already used
in Section~\ref{har:subsec:small}. The base field is the same; the
vanishing argument now uses the actual degree-zero direct image.

\begin{lemma}[Localized multiplier-ideal vanishing]\label{inj:local-nadel}
Let $Y$ be a smooth connected projective complex variety of dimension
$d>0$.  Suppose a line bundle $R$ has a singular Hermitian metric with
plurisubharmonic weights $\psi$.  Assume that on some coordinate ball
these weights are continuous and their curvature dominates a positive
smooth $(1,1)$-form.  Then
\[
 H^q\bigl(Y,\mathcal O_Y(K_Y+R)\otimes\mathcal I(\psi)\bigr)=0
 \qquad(q>0).
\]
\end{lemma}

\begin{proof}
We construct a globally strictly positive metric with exactly the same
multiplier ideal.

\smallskip\noindent\textit{Producing an integrable divisor metric.}
Choose a very ample $A$ with curvature $\omega_A$,
and a smooth metric $\psi_0$ on $R$.  For sufficiently small fixed
$\delta>0$ and arbitrarily large integers $l$, we first claim that
\begin{equation}\label{inj:jet-section}
 H^0\bigl(Y,\mathcal O_Y(lR-A)\otimes
       \mathcal I(l(1+\delta)\psi)\bigr)\ne0.
\end{equation}
On $lR$ use weights $l(1+\delta)\psi-l\delta\psi_0$, and add the
smoothly metrized line $j_lA-K_Y$, with
$j_l=\lceil C\delta l+C'\rceil$.  Fixed sufficiently large constants
$C,C'$ absorb the negative smooth curvatures and leave a positive
lower bound.  The multiplier ideal is
$\mathcal I_l=\mathcal I(l(1+\delta)\psi)$.
Within a smaller ball of strict positivity choose a point and a cutoff
logarithmic pole equal to $\log|z|^2$ near that point.  Add $\eta l$
times this pole, with $\eta>0$ fixed sufficiently small independently
of $\delta,l$.  Its negative curvature costs at most $O(\eta l)$
on the ball, where the original positive lower bound has order $l$.
Nadel vanishing \cite{Nadel1990}, in the strict-current form of
\cite[Theorem~5.11]{Demailly2011}, therefore gives surjectivity to jets
of order $\lfloor\eta l\rfloor-d$ at the point.  Continuity of
$\psi$ there makes $\mathcal I_l$ trivial near that point, and
polar-coordinate integration gives the ideal
$\mathfrak m^{\lfloor\eta l\rfloor-d+1}$ for the added pole.
Consequently
\[
 h^0\bigl(Y,(lR+j_lA)\otimes\mathcal I_l\bigr)\ge c l^d
\]
for some $c>0$ independent of sufficiently small $\delta$.

Remove $j_l+1$ copies of $A$.  At each step choose a smooth member of
$|A|$ containing no associated center of $\mathcal O_Y/\mathcal I_l$.
If a section vanishes on that member, division by its equation
preserves membership in $\mathcal I_l$, because that equation is a
nonzerodivisor on the quotient.  The number of conditions lost at a
step is at most $C''l^{d-1}$ uniformly for $0<\delta\le1$.
Indeed choose $p$ with $pA-R$ globally generated, restrict a general
nonzero section to the hyperplane, and bound the restriction space by
that of $(pl+h)A$ there, where $0\le h\le j_l$.  The Hilbert
polynomial and the hyperplane exact sequence give this bound.
The total loss is $O(\delta l^d)+O(l^{d-1})$.  Choosing $\delta$
small and then $l$ large proves \eqref{inj:jet-section}.

Let $\sigma$ be a nonzero section in \eqref{inj:jet-section}. In compatible
local frames, the weight $\phi_A$ of the positive metric on $A$ defines a divisor
metric on $R$ by
\[
 \psi'=\frac{1}{l}\bigl(\log|\sigma|^2+\phi_A\bigr).
\]
Its curvature dominates $\omega_A/l$. If
$V=\psi'-(1+\delta)\psi$, membership of $\sigma$ in the multiplier ideal
in \eqref{inj:jet-section} says exactly that $\exp(lV)$ is locally
integrable.

\smallskip\noindent\textit{Preserving the ideal.} Take
$l\ge1/\delta$ and put
$\rho=(1-\varepsilon)\psi+\varepsilon\psi'$.  The identity
\[
 -\psi=-\frac{\rho}{1+\varepsilon\delta}
             +\frac{\varepsilon V}{1+\varepsilon\delta}
\]
and H\"older's inequality, with the integrability of
$\exp(V/\delta)$, prove
$\mathcal I(\rho)\subseteq\mathcal I(\psi)$.
For the converse, strong openness
\cite[Theorem 1.1]{GuanZhou2015}, coherence, and compactness give
$\gamma>0$ with
$\mathcal I((1+\gamma)\psi)=\mathcal I(\psi)$ on a finite cover.
Small-exponent integrability of the divisor metric $\psi'$ gives
$\gamma'>0$ with $\exp(-\gamma'\psi')$ locally integrable.
Apply H\"older with
$s=(1-\varepsilon)/(1+\gamma)$ and $1-s$.  For small
$\varepsilon$ one has $\varepsilon/(1-s)<\gamma'$, proving
$\mathcal I(\psi)\subseteq\mathcal I(\rho)$.
The metric $\rho$ has a global positive curvature lower bound and
the same ideal.  Nadel vanishing proves the assertion.
\end{proof}

\begin{lemma}[Strictness from horizontal curvature]
\label{inj:strict-direct-image}
Let $p:V\to Y$ be a proper holomorphic submersion with compact fibers
and K\"ahler total space.  Let $D$ have a smooth semipositive
Hermitian metric.  Work on an open set where
$p_*(K_{V/Y}+D)$ is locally free and its fibers are the corresponding
section spaces.  Suppose the coefficient curvature dominates the
pullback of a positive base form on an open patch meeting $V_{y_0}$,
and let $u_0\in H^0(V_{y_0},K_{V_{y_0}}+D|_{V_{y_0}})$ be a
section that is nonzero on that patch. The curvature of the
direct-image metric paired with $u_0$ at $y_0$ is strictly positive
in every nonzero tangent direction of $Y$.
If a compact torus acts over the trivial base action and preserves the
metrics, let the entire invariant summand, or an entire character
summand, have rank one and suppose $u_0$ belongs to that summand.
Its quotient metric under the holomorphic orthogonal projection has
the same strictness at $y_0$.
\end{lemma}

\begin{proof}
We spell out the consequence of Berndtsson's curvature formula
\cite[Section 4, Proposition 4.2, Lemmas 4.3--4.4, and equation (4.8)]{Berndtsson2009},
since strict vertical positivity is not required.  Restrict to a disk
with coordinate $z$ in a chosen base direction, and let the fiber
dimension be $a$. If $a=0$, the direct image is locally the orthogonal
sum of the coefficient lines on the finitely many sheets. Its
curvature pairing is a sum of nonnegative terms, with a positive
term on the patch where $u_0$ is nonzero. Assume now $a>0$.
A holomorphic direct-image section extending $u_0$
with vanishing covariant derivative at zero has a smooth absolute
$(a,0)$ representative $u$ satisfying
\[
 D'u=dz\wedge\lambda,\qquad
 \bar\partial u=dz\wedge\eta,
 \qquad D'=\partial-\partial\phi\wedge.
\]
On the central fiber, $\lambda$ is orthogonal to holomorphic adjoint
forms and $\eta\wedge\omega$ is $\bar\partial$-exact.
Replace $u$ by $u-dz\wedge\xi$ so that $\lambda=0$ and
$\eta$ is primitive there.  These two requirements prescribe,
respectively, $\bar\partial^*(\omega\wedge\xi)$ and
$\bar\partial(\omega\wedge\xi)$, up to the K\"ahler-identity
constants.  Orthogonality and exactness place the two prescriptions
in their required images; Hodge decomposition solves them
simultaneously.  Lefschetz linear algebra identifies the relevant
$(a,1)$-forms with $\omega\wedge\xi$, for $\xi$ of type
$(a-1,0)$.

With $c_a=i^{a^2}$, the resulting curvature pairing is
\[
 c_a p_*\bigl(u\wedge\bar u\,e^{-\phi}
                   \wedge i\partial\bar\partial\phi\bigr)
 -i\,dz\wedge d\bar z\,c_a
       \int_{V_0}\eta\wedge\bar\eta\,e^{-\phi}.
\]
The second term is nonnegative by the primitive middle-degree
Hodge--Riemann sign. The first integrand is nonnegative everywhere.
On the prescribed patch the pullback lower bound contributes the
positive base form times the squared vertical norm of $u_0$.
This is positive on a nonempty open subset, so its integral is
strictly positive. For a torus action, averaging
is a holomorphic orthogonal projection.  The Chern connection
preserves the orthogonal character splitting, so the same calculation
applies on the invariant summand.
\end{proof}

\begin{proposition}[The exact ideal of an invariant adjoint line]
\label{inj:exact-adjoint-ideal}
Let $X,V,Y$ be smooth connected projective varieties with $\dim Y>0$,
let $\pi:V\to X$ be birational, and let $f:V\to Y$ be surjective.
Suppose an algebraic torus $T$ acts equivariantly on these maps and
trivially on $Y$.
Give $L=-K_X$ its natural linearization and a smooth semipositive
metric invariant under the maximal compact torus. Put $L'=\pi^*L$,
$J=K_{V/X}$, and let $t$ be its natural invariant Jacobian section.
Assume that
\[
 \mathcal B=(f_*\mathcal O_V(J))^T
\]
has rank one. Via division by $t$, view $\mathcal B$ as a sheaf of
rational functions on $Y$, and write
$\mathcal B^{\vee\vee}=\mathcal O_Y(B_0)$ and $R=B_0-K_Y$.
Then the invariant adjoint $L^2$ metric extends to a singular
semipositive metric $\psi_R$ on $R$, and
\begin{equation}\label{inj:exact-ideal}
 \mathcal O_Y(B_0)\otimes\mathcal I(\psi_R)=\mathcal B.
\end{equation}
If, on a nonempty patch in the smooth base-change locus where $t$ is
nonzero, the curvature of $L'$ dominates the pullback of a positive
base form, then $\psi_R$ is smooth and strictly positive on a base
coordinate ball.
\end{proposition}

\begin{proof}
The sheaf $\mathcal B$ is torsion-free. Its hull is determined by
regularity over the generic points of base primes. More precisely,
for a prime $P\subset Y$ its coefficient in $B_0$ is
\[
 b_P=\min_{Q\mapsto P}
          \left\lfloor\frac{\operatorname{ord}_Q(t)}{
          \operatorname{ord}_Q(f^*P)}\right\rfloor.
\]
These are exactly the regularity tests for the rational section
$(f^*a)t$. On the smooth base-change open, the entire invariant
adjoint summand is a line. Its fiber-integral metric is semipositive
by the compact-torus orthogonal quotient argument of
Proposition~\ref{har:prop:maximum}.

Put $k=\dim Y$. This metric extends plurisubharmonically across every base prime.
At a general point of a component attaining $b_P$, write the map
as $(z_1^e,z_2,\ldots,z_k)$, where $e>0$ is the pullback
multiplicity, and let $a$ be the order of $t$.
The relative canonical coefficient for the local frame of $R$ is
a nonvanishing smooth factor times
\[
 z_1^{a-b_Pe-(e-1)}.
\]
Since $0\le a-b_Pe<e$, this exponent is nonpositive.
Integration on one fixed small vertical polydisk bounds the squared
norm below by a positive constant near the base prime.  The weight
is therefore locally bounded above there.  Extend first along these
divisorial neighborhoods and then across codimension at least two;
write the extended weight as $\psi_R$.

Choose local frames $g$ of $\mathcal O_Y(B_0)$ and $\eta$ of $K_Y$.
The frame $g/\eta$ of $R$ corresponds on the smooth locus to the
relative adjoint section $(f^*g)t/f^*\eta$. For a holomorphic
function $a$, Fubini identifies integrability of $a$ with this frame's
squared norm, using the base volume from $\eta$, with square
integrability of the absolute top form $(f^*(ag))t$ with the smooth
$L'$ metric upstairs. The excluded analytic sets have measure
zero.  Such an integrable meromorphic top form has no divisorial
pole; smoothness upstairs then extends it across codimension two.
Conversely a holomorphic lift is integrable over inverse images of
relatively compact base charts by properness.  The lift is invariant
and the identification agrees on the generic fiber.  This proves
both inclusions in \eqref{inj:exact-ideal}, so it retains any
codimension-two difference between $\mathcal B$ and its hull.

For the final assertion apply Lemma~\ref{inj:strict-direct-image}
to the invariant section determined by $t$ on the stated patch.
The fiber-integral metric is smooth on the base-change open, so
its strictness at one point gives a positive lower curvature bound
on a smaller coordinate ball.
\end{proof}

\begin{proposition}[Invariant sections by a degree-zero field]
\label{inj:degree-zero-route}
Let $X$ be smooth connected projective with
$H^q(X,\mathcal O_X)=0$ for $q>0$.  Suppose $L=-K_X$ has a smooth
semipositive metric and an algebraic torus $T$ acts on $X$.
Then $H^0(X,mL)^T\ne0$ for some $m>0$, with the natural
anticanonical linearization.
\end{proposition}

\begin{proof}
\smallskip\noindent\textit{The null-pole field and the relative section.}
Average a metric on $L$ and a K\"ahler form $\omega$ over the maximal
compact torus, and write $\alpha\ge0$ for the resulting curvature
form, in first-Chern normalization.  Proposition~\ref{har:prop:index}, with its actual value at zero,
followed by the equivariant hard Lefschetz construction in
Lemma~\ref{har:lem:forms},
has the following consequence: if the asserted invariant sections do
not exist, there is a fixed $q>0$ and unbounded positive $j$ with
\[
 0\ne u_j\in H^0(X,\Omega_X^{n-q}\otimes jL)^T.
\]
The coefficient in hard Lefschetz is $j=m+1$, since
$K_X+jL=mL$ equivariantly.  In the Bochner identity for the associated
harmonic $(n,q)$-forms the curvature term is a sum of $q$ eigenvalues
of $j\alpha$.  Full rank on an open set would force vanishing there,
hence vanishing of the holomorphic Lefschetz preimage everywhere.
Thus $r=\max\operatorname{rank}\alpha<n$.

Define the degree of a divisor by
\[
 \ell(D)=[D]\,[\alpha]^r[\omega]^{n-r-1}.
\]
It is nonnegative on effective divisors and vanishes on $L$.
Lemma~\ref{har:lem:forms} gives a linearized determinant line $S$
and, for unbounded positive integers $p$, nonzero invariant sections
\[
 s_p\in H^0(X,S+pL)^T.
\]
Put $N_p=\operatorname{div}(s_p)$. Lemma~\ref{har:lem:null-divisors}
gives $\ell(N_p)=0$: prescribed Ricci curvature first gives
$\ell(S)\le0$, and effectivity supplies the reverse inequality.
The degree-zero divisors are exactly the curvature-null divisors
of Section~\ref{har:subsec:small}. We use their invariant pole field
again, now to construct the exact direct-image ideal.

Let $K$ be the field of invariant rational functions whose polar
divisors have degree zero.  Sums and products preserve this property;
inverses do also, because polar and zero divisors are linearly
equivalent and effective degrees are nonnegative.  The field is
relatively algebraically closed in $\mathbb C(X)$.  Indeed a root
of a monic equation over $K$ can have poles only on the polar
divisors of its coefficients, by the discrete-valuation inequality.
The connected torus fixes each member of the finite root set, so the
root is invariant as well.  The field is finitely generated: choose
a transcendence basis over $\mathbb C$; algebraic subextensions
over its rational field inside the finitely generated ambient field
have bounded degrees, after adjoining an ambient separating
transcendence basis.  Their compositum is therefore finite.

Take a smooth projective model $Y$ of $K$.  Flatten the graph of the
induced rational map, and normalize its main component, to obtain a
normal equidimensional $G\to Y$ birational to $X$; then take an
equivariant resolution $V\to G$.  The action on $Y$ is trivial. Write $f_G:G\to Y$ for the normal-model morphism and $L_G$ for the pullback of $L$ to $G$.
Write $\pi:V\to X$, $f:V\to Y$, $L'=\pi^*L$, and let
$t$ be the invariant Jacobian section of
$K_V+L'=\mathcal O_V(J)$, with $J=K_{V/X}\ge0$.
The generic fiber is geometrically integral because $K$ is relatively
algebraically closed.  These constructions use only equivariant
resolution of the graph and the Hilbert-family flattening; they do
not assert equidimensionality of the smooth resolution.

Ratios of products of the $s_p$ having the same bidegree in $(S,L)$
belong to $K$.  Their horizontal divisor on $G$ is zero, so for
$p>p_1>p_0$,
\[
 (p-p_0)N_{p_1}^{\rm hor}
 =(p-p_1)N_{p_0}^{\rm hor}+(p_1-p_0)N_p^{\rm hor}.
\]
Equivalently,
\[
 N_p^{\rm hor}=N_{p_0}^{\rm hor}
 +\frac{p-p_0}{p_1-p_0}
       (N_{p_1}^{\rm hor}-N_{p_0}^{\rm hor}).
\]
Effectivity for unbounded $p$ forces
$N_{p_1}^{\rm hor}-N_{p_0}^{\rm hor}\ge0$.
Put $d_0=p_1-p_0>0$ and
$\sigma_0=s_{p_1}/s_{p_0}$, a rational invariant section of
$d_0L$ with effective horizontal divisor. We use the same symbol
for its pullback to $G$. If $Y$ is a point, it has no poles and
proves the assertion. Assume $\dim Y=k>0$.

For each prime $P\subset Y$, put
\[
 c_P=\min_{Q\mapsto P}
       \frac{\operatorname{ord}_Q(\sigma_0)}
            {\operatorname{ord}_Q(f_G^*P)}.
\]
Only finitely many $c_P$ are nonzero. Choose a positive integer $h$
clearing their denominators, and set
\[
 d=hd_0,\qquad \sigma=\sigma_0^h,\qquad
 B=\sum_P h c_P P.
\]
Thus $\sigma$ is a rational invariant section of $dL_G$, and $B$
is an integral divisor with trivial action. Divisorial normalization
gives
\begin{equation}\label{inj:invariant-base-decomposition}
 s_D=\sigma(f_G^*1_B)^{-1}
     \in H^0(G,dL_G-f_G^*B)^T,
 \qquad D=\operatorname{div}(s_D)\ge0.
\end{equation}
Here $1_B$ is the rational canonical section of $\mathcal O_Y(B)$.
At least one component above every base prime has coefficient zero
in $D$. We also write $\sigma$ for its pullback to $V$.

\smallskip\noindent\textit{Adjoint rigidity and a horizontal curvature bound.}
For every $m\ge0$,
\begin{equation}\label{inj:degree-rank}
 \operatorname{rank}
 \bigl(f_*\mathcal O_V(J+mdL')\bigr)^T=1.
\end{equation}
There is a generic section $t\sigma^m$.  If the invariant rank
were larger, a large very ample base twist would give two invariant
global sections independent over $K$.  Their pushed divisors on $X$
have degree zero.  For the base twist this follows by taking the
ratio of two general base hyperplanes: their pullbacks to the
equidimensional $G$ have no common prime component, so their pushed
divisors are the actual polar and zero divisors of an element of
$K$.  The remaining class is $mdL$, and the exceptional Jacobian
pushes to zero.  Thus the ratio of the two sections would again
belong to $K$, contradicting their independence.  This proves
\eqref{inj:degree-rank}, including $m=0$.

At a rank-$r$ point of $\alpha$, the differential of every rational
function in $K$ kills $\ker\alpha$ wherever it is regular.
For a regular value, the closure of its level hypersurface has
degree zero.  The nonnegative form
$\alpha^r\wedge\omega^{n-r-1}$ therefore vanishes on that
hypersurface.  A tangent hyperplane failing to contain $\ker\alpha$
would still carry rank $r$, a contradiction.  Apply this to $k$
independent base coordinates. It gives $k\le r$. Intersect the
maximal-rank open with the dense smooth base-change locus and the
complement of the exceptional locus. On a small patch there,
\begin{equation}\label{inj:horizontal-patch}
 \pi^*\alpha\ge c f^*\omega_Y\qquad(c>0).
\end{equation}
This is semipositive linear algebra on a complement of the kernel,
not an inference from numerical degree alone.  If $r=0$, a nonzero
effective divisor cannot have degree zero, so $k>0$ cannot occur.

Apply Proposition~\ref{har:prop:maximum} to the equivariant normal equidimensional model $G\to Y$, with $P=L$ naturally linearized and the rational section $\sigma$. Equation~\eqref{inj:degree-rank} verifies the invariant adjoint ranks, and \eqref{inj:invariant-base-decomposition} is its vertical normalization. The result is a positive continuous semipositive metric on $B$.

\smallskip\noindent\textit{The exact ideal and interpolation from zero.}
The varying twist is now metrized. The degree-zero metric must retain
the entire direct image, including its codimension-two defect.

Apply Proposition~\ref{inj:exact-adjoint-ideal} to
$\mathcal B=(f_*\mathcal O_V(J))^T$. Equation~\eqref{inj:degree-rank}
at $m=0$ gives its rank, and the patch in
\eqref{inj:horizontal-patch} is disjoint from the exceptional locus,
so $t$ is nonzero there. The proposition supplies the metric
$\psi_R$ on $R=B_0-K_Y$, strictly positive on a base ball, with the
exact identity \eqref{inj:exact-ideal}.

Adding $m$ times the continuous bounded metric on $B$ preserves the
ideal and the strictness on the ball.  Lemma~\ref{inj:local-nadel}
therefore gives
\[
 H^q(Y,\mathcal B\otimes\mathcal O_Y(mB))=0
 \qquad(q>0,\ m\ge0).
\]
At zero, canonical pushforward gives
$h^0(Y,\mathcal B)=h^0(V,\mathcal O_V(J))^T=1$.
Thus the Euler polynomial has value one at zero.  For some $m>0$ lift a nonzero section
and multiply by $s_D^m$.  The result is an invariant section of
$J+mdL'$.  Remove the exceptional Jacobian divisor by pushforward
to the smooth $X$.  It gives the required natural invariant section
of $mdL$.
\end{proof}
